% Generated by `taskgen course examples` from ecc.rho (seed 12).
% Do not edit: change examples.yaml and regenerate.
\begin{example}[$\rho$-метод, необратимый коэффициент]\label{ex:ecdlp-rho-composite}
На кривой $E\colon y^2 = x^3 + 27x + 46$ над $\mathbb{F}_{47}$ точка $P = (2, 25)$ имеет порядок $N = 14$. Решить ECDLP $Q = nP$, $Q = (10, 0)$, $\rho$-методом Полларда: $R_{j+1} = R_j + P$, $2R_j$ или $R_j + Q$, если $x(R_j) \bmod 3$ равно $0$, $1$ или $2$ соответственно (точку $\mathcal{O}$ относим к первому множеству); $R_0 = P$; коллизию искать методом Флойда ($R_i = R_{2i}$).

\par\noindent\textit{Решение.} Храним $R_j = a_j P + b_j Q$; при $R \gets R + P$: $a \gets a + 1$; при $R \gets 2R$: $(a, b) \gets (2a, 2b)$; при $R \gets R + Q$: $b \gets b + 1$ (всё по модулю $N = 14$). Формулы: $\lambda = \frac{3x_1^2 + a}{2y_1}$ (удвоение), $\lambda = \frac{y_2 - y_1}{x_2 - x_1}$ (сложение), $x_3 = \lambda^2 - x_1 - x_2$, $y_3 = \lambda(x_1 - x_3) - y_1$; если $x_1 = x_2$, $y_1 = -y_2$, то сумма равна $\mathcal{O}$; $\mathcal{O}$ "--- нейтральный элемент. Координаты по модулю $47$, $a_E = 27$.\par Последовательность $R_j$ ($s = x(R_j) \bmod 3$) до $j = 2i = 12$:\[\begin{array}{r|c|c|r|r}j & R_j & s & a_j & b_j \\ \hline 0 & (2, 25) & 2 & 1 & 0 \\ 1 & (22, 14) & 1 & 1 & 1 \\ 2 & (24, 18) & 0 & 2 & 2 \\ 3 & (1, 11) & 1 & 3 & 2 \\ 4 & (22, 33) & 1 & 6 & 4 \\ 5 & (24, 29) & 0 & 12 & 8 \\ 6 & (2, 22) & 2 & 13 & 8 \\ 7 & (22, 33) & 1 & 13 & 9 \\ 8 & (24, 29) & 0 & 12 & 4 \\ 9 & (2, 22) & 2 & 13 & 4 \\ 10 & (22, 33) & 1 & 13 & 5 \\ 11 & (24, 29) & 0 & 12 & 10 \\ 12 & (2, 22) & 2 & 13 & 10 \\ \end{array}\]Флойд: $R_{1} \ne R_{2}$, $R_{2} \ne R_{4}$, $R_{3} \ne R_{6}$, $R_{4} \ne R_{8}$, $R_{5} \ne R_{10}$, $R_{6} = R_{12} = (2, 22)$. Отсюда $13 P + 8 Q = 13 P + 10 Q$, т.\,е. $(b_i - b_{2i})\, n \equiv a_{2i} - a_i$: \[12\, n \equiv 0 \pmod{14}.\]\par $d = \gcd(12, 14) = 2$ (коэффициент необратим, $N$ составное); $d \mid 0$, делим на $d$: $6\, n \equiv 0 \pmod{7}$, $n \equiv 0 \pmod{7}$. Кандидаты $n = 0 + k \cdot 7$, $0 \le k < 2$; проверяем: $0P = \mathcal{O}$; $7P = (10, 0)$ $= Q$. Проверка: $7P = Q$.\[n = 7\ (\textrm{mod}\ 14).\]
\end{example}
